\bookchapter{Capitolo 35: Bambini e telefoni — guida all'educazione nell'era digitale}

\chapterintro{Questo capitolo non racconta che “il telefono è un mostro”, né sponsors “chi prima comincia prima diventa intelligente”. Procediamo per fasce d'età: device, regole, alfabetizzazione, salute — un passo alla volta.}

\booksection{A quale età dare il telefono al bambino?}

Non c'è una risposta standard, ma c'è una baseline di età. Se lo dai troppo presto, la corteccia prefrontale del bambino non è ancora matura e l'autocontrollo non regge il “design persuasivo” degli algoritmi; troppo tardi, e rischia l'“emarginazione sociale” tra i pari. Ecco i consigli per fascia d'età:

\begin{itemize}
  \item \textbf{3-6 anni}: possibilmente non tocca il telefono. Se deve usarlo, non più di 15 minuti per volta, con contenuti prevalentemente cartoni educativi selezionati dai genitori, e\textbf{sempre guardando insieme}. In questa fascia lo schermo ha l'impatto maggiore su vista e attenzione.
  \item \textbf{7-9 anni}: può "prendere in prestito" il dispositivo dei genitori per lezioni online, ascoltare storie, cercare informazioni. Inizia a costruire la regola "chiedi prima di usare, restituisci dopo", senza configurare un telefono dedicato.
  \item \textbf{10-12 anni}: si può considerare un "feature phone" o un orologio per bambini, per il contatto andata-ritorno da scuola. Se serve davvero uno smartphone, attiva controlli parentali rigorosi.
  \item \textbf{oltre 13 anni}: entra alle medie, le esigenze sociali crescono, puoi dare uno smartphone, ma accompagnato da un "patto d'uso" e da verifiche periodiche.
\end{itemize}

\begin{tcbinfo}{Principio "abbasta e avanza"}
Il primo telefono del bambino non deve essere un flagship. Le “prestazioni eccedenti” del processore spesso finiscono per far girare giochi. Scegliere un telefono di fascia media con buona batteria, schermo riposante e prezzo sui 1000 yuan, oppure un ricondizionato ufficiale, agevola il controllo. Budget di riferimento: 1500-2500 yuan bastano per la stragrande maggioranza degli studenti.

\end{tcbinfo}

\booksubsection{L'orologio intelligente come soluzione transitoria}

Se vuoi solo risolvere il problema “riesci a contattarlo dopo scuola”, un orologio per bambini è più adatto del telefono. Può solo chiamare, inviare messaggi vocali, localizzare — niente browser né brevi video, e limita naturalmente le “porte di dipendenza”:

\begin{itemize}
  \item \textbf{Orologio telefonico XiaoTianCai}: l'ecosistema più completo; il "touch per aggiungere amici" tra coetanei è il suo punto di forza, ma l'ecosistema app dedicato è piuttosto chiuso.
  \item \textbf{Orologio per bambini Huawei}: buona integrazione col telefono Huawei, localizzazione precisa, supporta videotelefonate, ottimo rapporto qualità-prezzo.
  \item \textbf{Orologio per bambini Xiaomi}: prezzo accessibile, funzioni di base complete, adatto a famiglie con budget limitato.
\end{itemize}

Quando compri l'orologio, attenzione: scegli modelli che supportano\textbf{4G full-network}(2G/3G sono in fase di dismissione), e verifica le tariffe del piano bambini dell'operatore.

\booksection{Guida completa ai controlli parentali}

Prima di consegnare il dispositivo al bambino, la prima cosa non è accenderlo, ma configurare i controlli parentali. Entrambe le piattaforme hanno soluzioni native:

\booksubsection{iOS: Tempo di utilizzo}

\begin{itemize}
  \item \textbf{Limiti app}: Impostazioni > Tempo di utilizzo > Limiti app, puoi definire il tempo massimo giornaliero per categoria di app (ad esempio 30 minuti per i giochi).
  \item \textbf{Tempo di inattività}: imposta una "fase di spegnimento" quotidiana (ad esempio 22:00-7:00), in cui tutte le app si disattivano, resta solo il telefono.
  \item \textbf{Limiti di contenuto}: filtra film, pagine web e installazione di app per fascia d'età, e disabilita gli acquisti in-app.
  \item \textbf{In famiglia}: gestisci da remoto il dispositivo del bambino dall'iPhone del genitore, approva le richieste di download e consulta i report d'uso.
\end{itemize}

\booksubsection{Android: Google Family Link / controlli parentali di marca}

\begin{itemize}
  \item \textbf{Google Family Link}: soluzione ufficiale Google, gestisce da remoto il tempo d'uso del dispositivo, lo blocca e mostra i report delle attività. Sugli Android cinesi dipende dalla disponibilità per modello.
  \item \textbf{Spazio bambini Xiaomi}: Impostazioni > Spazio bambini, consente di impostare tempo d'uso, app in whitelist, l'uscita richiede la password del genitore.
  \item \textbf{Modalità studente Huawei}: Impostazioni > Uso sano del telefono > Modalità studente, limita fasce orarie e app, supporta la gestione remota.
  \item \textbf{Modalità bambini OPPO/vivo}: si attiva in "Impostazioni" o nell'app "Spazio bambini" dedicata, funzionamento simile.
\end{itemize}

\begin{center}
\begin{tabularx}{\linewidth}{|X|X|X|}
\hline
\textbf{Funzione} & \textbf{iOS Tempo di utilizzo} & \textbf{Android Family Link / soluzioni di marca} \tabularnewline
\hline
Limiti app & sì, per categoria & sì, per app \tabularnewline
Tempo di inattività & sì & sì (con differenze tra marchi) \tabularnewline
Filtro contenuti & a livello di sistema, più radicale & dipende dalla marca, disomogeneo \tabularnewline
Gestione remota & In famiglia & Family Link / cloud di marca \tabularnewline
Approvazione installazione app & richiede approvazione del genitore & supportato da Family Link \tabularnewline
\hline
\end{tabularx}
\end{center}

\begin{tcbwarn}{Il controllo parentale non è "una volta per tutte"}
I bambini trovano sempre il modo di aggirarlo (cercando online “come sbloccare il tempo di utilizzo” si trova una marea di tutorial). I mezzi tecnici sono solo una baseline; la vera difesa è la comunicazione e la fiducia genitore-figlio. Si consiglia di rivedere insieme il report d'uso ogni mese, invece di monitorare unilateralmente.

\end{tcbwarn}

\booksection{Modalità adolescente}

Le principali piattaforme di contenuti hanno una “modalità adolescente” integrata: una volta attivata limita il tempo d'uso, filtra i contenuti inadatti e disabilita ricariche e mance. È una funzione obbligatoria richiesta dalla legislazione nazionale, ma\textbf{deve essere attivata dal genitore, che deve impostare la password}:

\begin{itemize}
  \item \textbf{Modalità adolescente di Douyin/Kuaishou}: 40 minuti al giorno (disabilitata 22:00-6:00), mostra solo video selezionati di divulgazione scientifica ed educativi, senza live streaming, mance o ricariche.
  \item \textbf{Modalità adolescente di WeChat}: limita WeChat Pay (richiede autorizzazione del genitore con tetto), limita i contenuti degli account pubblici e dei canali video; Scuotiti, Persone vicine e altre funzioni sono disabilitate per default.
  \item \textbf{Sistema anti-dipendenza dei giochi}: autenticazione real-name + limiti di tempo a livello nazionale. I minori possono accedere ai giochi solo venerdì, sabato, domenica e festività, dalle 20:00 alle 21:00, 1 ora al giorno.
\end{itemize}

\textbf{Non usare la carta d'identità del genitore per il “prestito account”}\par

Molti bambini registrano i conti gioco con la carta d'identità dei genitori per aggirare l'anti-dipendenza. Questo viola il contratto utente e fa perdere al bambino la protezione. I genitori dovrebbero custodire le proprie password di pagamento e i propri dati identificativi e non “prestare” account al bambino per comodità.

\booksection{Educazione all'alfabetizzazione digitale}

La tecnologia intercetta il 90\% dei rischi; il restante 10\% dipende dal giudizio del bambino. Più che “non navigare”, conta insegnargli “come navigare”:

\begin{itemize}
  \item \textbf{Protezione della privacy}: non rivelare a estranei nome vero, scuola, indirizzo di casa; non inviare foto personali, documento d'identità, foto di famiglia; la localizzazione è disattivata per default, attivata solo quando serve.
  \item \textbf{Riconoscimento e gestione del cyberbullismo}: insulti, isolamento, fotomontaggi offensivi sono tutti bullismo. Dì al bambino: non rispondere con insulti, non cancellare le prove, avverti subito i genitori. Su WeChat puoi tenere premuto sul messaggio per segnalarlo e la piattaforma tratterà l'account dell'aggressore.
  \item \textbf{Spirito critico}: insegnagli a distinguere "fatti" da "opinioni", a riconoscere clickbait e titoli shock; a capire che l'algoritmo spinge solo ciò che vuoi vedere, non rappresenta il mondo intero; di fronte ad annunci "una torta dal cielo", non credere mai.
  \item \textbf{Sicurezza online}: non cliccare link sconosciuti, non scansionare QR code di origine dubbia, non scaricare app da store non ufficiali; non condividere la password con nessuno (nemmeno con il "migliore amico"); chiunque chieda un codice di verifica va considerato un truffatore.
\end{itemize}

\begin{tcbpractice}{Fai: firma con tuo figlio un "Patto di utilizzo del telefono"}
Scrivere le regole, firmarle entrambi e affiggerle in vista è molto più efficace di un accordo verbale. Si suggerisce di includere:

\begin{itemize}
  \item Fasce orarie: ad esempio "dopo i compiti, 19:30-20:30 disponibile".
  \item Zone vietate: a tavola, in camera prima di dormire, in classe il telefono va ritirato.
  \item Confini di contenuto: non installare app non comunicate ai genitori, non chattare in privato con estranei.
  \item Conseguenze in caso di violazione: ad esempio "una volta in ritardo, il tempo di utilizzo del giorno dopo si dimezza".
  \item Impegno del genitore: non "scorrere il telefono" davanti al bambino oltre il tempo pattuito, dare l'esempio.
\end{itemize}

\end{tcbpractice}

\booksection{Abitudini di uso sano}

Occhi, collo e sonno del bambino sono in fase di sviluppo e il telefono li danneggia più di quanto faccia con gli adulti:

\begin{itemize}
  \item \textbf{Tempo schermo consigliato}: 3-6 anni non più di 1 ora al giorno; 6-12 anni niente limiti per l'apprendimento, intrattenimento non oltre 1 ora; sopra i 12 anni da concordare col bambino, ma intrattenimento non oltre 2 ore.
  \item \textbf{Abitudini per la vista}: segui la "regola 20-20-20" — ogni 20 minuti di schermo, guarda a 20 piedi (circa 6 metri) di distanza per 20 secondi; niente telefono nell'ora prima di dormire, la luce blu inibisce la melatonina e altera il sonno.
  \item \textbf{Gestione della postura}: schermo all'altezza degli occhi, non a testa bassa o spalle curve; non più di 30 minuti consecutivi, poi alzati e muoviti.
  \item \textbf{Regole condivise in famiglia}: tavolo, camera da letto, studio come "zone senza telefono"; quando si guarda un film insieme o si fa attività all'aperto, i telefoni si ritirano: "mettere via il telefono" diventi cultura di famiglia, non punizione per il bambino.
\end{itemize}

\begin{tcbinfo}{Il genitore è il miglior "firewall"}
Le ricerche mostrano che il tempo di utilizzo del telefono da parte dei genitori è fortemente correlato a quello dei figli. Pretendere che il bambino non guardi il telefono mentre tu scori brevi video fino a tarda notte rende le regole inefficaci. Per cambiare il bambino, comincia dal mettere giù il tuo telefono.

\end{tcbinfo}

\booksection{Eredità digitale e minori}

Questa sezione è spesso trascurata, ma importante. I beni digitali del bambino (account, foto, saldi ricaricati, attrezzatura di gioco) affrontano anch'essi problemi di successione e privacy:

\begin{itemize}
  \item \textbf{Successione dei beni digitali dei minori}: in caso di premorienza del bambino, la successione di account social, account di gioco e foto su cloud segue lo stesso iter degli adulti, con certificato di morte e certificato di parentela. Si consiglia ai genitori di configurare in anticipo un "contatto per l'eredità" o "contatto di emergenza" per gli account importanti del bambino (ad esempio il contatto per l'eredità Apple, la successione dell'account WeChat).
  \item \textbf{Privacy digitale dei minori}: la "Legge sulla protezione delle informazioni personali" stabilisce che il trattamento dei dati personali di minori sotto i 14 anni richiede il consenso del tutore e deve seguire regole di privacy dedicate. I genitori dovrebbero verificare periodicamente le politiche di privacy delle app più usate dal bambino e disattivare le opzioni "raccomandazioni personalizzate" e "condivisione dati".
  \item \textbf{Consigli per la registrazione degli account}: gli account registrati per il bambino dovrebbero usare email/numero di telefono controllabili dal genitore, per una gestione unificata durante la minore età; a maggiore età si potrà fare il passaggio dell'account.
\end{itemize}

\begin{tcbwarn}{Attenzione all'altra faccia dell'"eredità digitale dei bambini"}
Il bambino viene “registrato” fin dalla nascita — foto di crescita pubblicate dai genitori, video inviati dall'asilo, dati di salute dell'orologio intelligente. Una volta caricati, questi dati sono difficili da eliminare del tutto. Prima di pubblicare, pensaci: tra dieci anni il bambino si vergognerà di questa foto? Rispettare la personalità digitale del bambino comincia dal trattenersi dal “postare il bambino”.

\end{tcbwarn}

\begin{tcbsummary}{Riepilogo del capitolo}
\begin{itemize}
  \item Non c'è risposta standard sull'età, ma 3-6 anni possibilmente senza, prima dei 13 usare l'orologio per bambini o un feature phone come transizione
  \item "Abasta e avanza": il primo telefono sceltilo di fascia media, non flagship; prestazioni eccedenti = tentazioni eccedenti
  \item Il Tempo di utilizzo di iOS e Family Link/controlli di marca Android sono due difese tecniche, meglio combinate
  \item La modalità adolescente e l'anti-dipendenza dei giochi sono obbligo di legge, ma vanno attivate e protette col password dal genitore
  \item L'alfabetizzazione digitale conta più dei filtri tecnici: insegnare al bambino a riconoscere bullismo, proteggere la privacy, leggere criticamente i contenuti
  \item Le abitudini sane cominciano dal genitore — tu metti giù il telefono, e il bambino farà lo stesso
  \item Eredità digitale e privacy dei minori sono temi emergenti: configura fin da ora il "contatto per l'eredità" e trattieni il postare il bambino
\end{itemize}

\end{tcbsummary}
